%
% File:    ./notes/logic/ch01-intro.tex
% Author:  Jiří Kučera <sanczes AT gmail.com>
% Date:    2026-10-21 00:46:06 +0200
% Project: Math Notes
% Brief:   Notes about mathematical logic, chapter 1
%
% SPDX-License-Identifier: CC-BY-SA-4.0
%

\chapter Introduction (intro).

\refsecformat \NameTitle

This document serves as a collection of notes taken from various sources about
mathematical logic, mostly from \cite{Men15}. The content of this document is
organized into three\footnote{A simple foot note.} chapters.

\refsec{ch:intro}, first outlines the structure of this document and then, in
\refsec{sec:funterm}, it gives a fundamental terminology and notation used in
the following chapters.

\refsec{ch:pcalc}, is dedicated to the propositional calculus.

Finally, \refsec{ch:folog}, is dedicated to the first-order logic.

\section Fundamental Terminology (funterm).

Mathematical logic is not a small area. Actually, it spans across many other
fields of mathematics, like the computer science theory or set theory. To
demonstrate it, lets take a look on axiomatic set theory. Axiomatic set theory
precisely defines a set via axioms written in mathematical logic. On the other
hand, mathematical logic uses notions of set and function to describe its core
concepts, as it will be shown later. Thus, these two areas of mathematics are
tightly interconnected.
